\section{Generazione casuale del viaggio}
\vspace{-5pt}
Dato il nome di una città, un tempo disponibile e un utente, l'algoritmo genera un viaggio casuale composto da un insieme di monumenti visitabili entro il budget temporale assegnato. 
La costruzione del viaggio tiene conto sia dei tempi di visita dei monumenti sia dei tempi di spostamento stimati tra la posizione corrente e la destinazione, oltre al tempo necessario per rientrare al punto di partenza.

L'algoritmo implementato combina diverse fasi principali:
\vspace{-5pt}
\begin{enumerate}
    \setlength{\itemsep}{0pt}
    \setlength{\parskip}{0pt}
    \item \textbf{Recupero dei dati iniziali}: viene identificato l'utente e vengono caricati tutti i monumenti della città selezionata.
    \item \textbf{Geocodifica del punto di partenza}: la città viene trasformata in coordinate geografiche.
    \item \textbf{Filtraggio dei monumenti}: vengono considerati solo i monumenti con coordinate valide e distanti al massimo 10 km dal punto di partenza.
    \item \textbf{Selezione casuale vincolata}: i monumenti vengono mescolati casualmente e aggiunti al viaggio solo se il tempo residuo consente sia la visita sia il ritorno al punto iniziale.
\end{enumerate}

\vspace{-10pt}
\subsection*{Pseudocodice}
\vspace{-5pt}

\captionof{algorithm}{Generazione di un viaggio casuale vincolato dal tempo}
\begin{algorithmic}[1]
\Function{generateRandomTrip}{city, availableMinutes, username}
    \State user $\gets$ \Call{findUserByUsername}{username}
    \State monuments $\gets$ \Call{getMonumentsByCity}{city}
    \If{monuments is empty}
        \State \textbf{error} ``Nessun monumento trovato per la città''
    \EndIf

    \State startPoint $\gets$ city
    \State (startLat, startLon) $\gets$ \Call{geocode}{startPoint}
    \If{startLat = 0 \textbf{and} startLon = 0}
        \State \textbf{error} ``Impossibile trovare le coordinate della città''
    \EndIf

    \State monuments $\gets$ \Call{filterValidAndNearMonuments}{monuments, startLat, startLon, 10km}
    \If{monuments is empty}
        \State \textbf{error} ``Nessun monumento trovato entro il raggio consentito''
    \EndIf

    \State \Call{shuffle}{monuments}
    \State remaining $\gets$ monuments
    \State stages $\gets$ lista vuota
    \State usedSeconds $\gets$ 0
    \State budgetSeconds $\gets$ availableMinutes $\cdot$ 60
    \State currentLat $\gets$ startLat, currentLon $\gets$ startLon

    \While{remaining is not empty \textbf{and} size(stages) $<$ 10}
        \State chosen $\gets$ null
        \For{each candidate in remaining}
            \State visitMin $\gets$ \Call{estimateVisitMinutes}{candidate}
            \State travelToSec $\gets$ \Call{estimateTravelSeconds}{currentLat, currentLon, candidate.lat, candidate.lon}
            \State returnSec $\gets$ \Call{estimateTravelSeconds}{candidate.lat, candidate.lon, startLat, startLon}
            \State needed $\gets$ travelToSec + visitMin $\cdot 60 + returnSec$
            \If{usedSeconds + needed $\leq$ budgetSeconds}
                \State chosen $\gets$ candidate
                \State usedSeconds $\gets$ usedSeconds + travelToSec + visitMin $\cdot 60$
                \State \textbf{break}
            \EndIf
        \EndFor

        \If{chosen = null}
            \State \textbf{break}
        \EndIf

        \State aggiungi a stages il monumento scelto con durata di visita visitMin
        \State currentLat $\gets$ chosen.lat, currentLon $\gets$ chosen.lon
        \State rimuovi chosen da remaining
    \EndWhile

    \If{stages is empty}
        \State \textbf{error} ``Tempo disponibile insufficiente per visitare almeno un monumento''
    \EndIf

    \State trip $\gets$ crea viaggio con utente, città, punto di partenza e stages
    \State \Return \Call{saveTrip}{trip}
\EndFunction
\end{algorithmic}



\subsection*{Analisi della complessità}
\vspace{-5pt}
L'algoritmo parte recuperando i monumenti associati alla città e verificando la validità del punto di partenza tramite geocodifica. 
La fase di filtraggio dei monumenti, che elimina quelli privi di coordinate o troppo lontani dal punto iniziale, richiede una scansione dell'insieme dei monumenti e ha quindi complessità temporale $O(n)$, dove $n$ è il numero totale di monumenti disponibili.

Successivamente, la lista dei monumenti viene mescolata casualmente, operazione che ha complessità $O(n)$.
La fase principale è il ciclo di selezione dei monumenti: per ogni tappa del viaggio, l'algoritmo esamina i monumenti rimanenti fino a trovarne uno compatibile con il tempo residuo. 
Nel caso teorico senza vincolo sul numero massimo di tappe, questa procedura può portare a una complessità fino a $O(n^2)$, poiché per ogni monumento selezionato viene eseguita una scansione della lista dei candidati rimanenti.

Tuttavia, nel codice implementato il numero di tappe è limitato a 10, quindi il costo della fase principale è in pratica proporzionale a $10n$, che si può considerare $O(n)$ rispetto al numero di monumenti. 
Di conseguenza, la complessità complessiva dell'algoritmo è dominata dalle scansioni dei monumenti e risulta lineare in pratica, cioè $O(n)$, mentre una stima più generale senza il limite fisso sulle tappe porta a $O(n^2)$.